\chapter{Najkrótsze drogi ważone: Dijkstra, Bellman--Ford, A* i hierarchie kontrakcji}
\label{ch:tg-shortest-paths}

W rozdziale \ref{ch:tg-bfs-dfs} algorytm BFS wyznacza
najkrótsze drogi ze względu na liczbę krawędzi.
Teraz kosztem drogi będzie suma jej wag.
Poznamy algorytmy Dijkstry, Bellmana--Forda,
Floyda--Warshalla i A*, wyszukiwanie dwukierunkowe
oraz hierarchie kontrakcji (ang. \emph{Contraction Hierarchies}, CH).
Nie ma algorytmu najszybszego na wszystkich grafach
i przy wszystkich rodzajach zapytań.

\section{Grafy ważone, odległość i relaksacja}
\label{sec:tg-weighted-data}

Niech $D=(V,A)$ będzie skończonym grafem skierowanym
bez pętli, $n=|V|$, $m=|A|$ i niech
$w\colon A\to\mathbb R$ będzie funkcją wag.
Wagę spaceru $P=(v_0,\ldots,v_k)$ definiujemy jako
$w(P)=\sum_{i=0}^{k-1}w(v_i,v_{i+1})$.
Graf nieskierowany modelujemy dwoma przeciwnymi
łukami tej samej wagi.

\begin{definicja}[Odległość ważona]
\label{def:tg-weighted-distance}
Dla $s,t\in V$ niech $\delta(s,t)$ oznacza infimum
wag skierowanych spacerów z $s$ do $t$.
Gdy nie ma takiego spaceru, przyjmujemy $+\infty$.
Gdy istnieje osiągalny z $s$ cykl o ujemnej wadze,
z którego można dojść do $t$, przyjmujemy $-\infty$.
\end{definicja}

\begin{lemat}[Realizacja odległości]
\label{lem:tg-weighted-realization}
Jeżeli $t$ jest osiągalny z $s$ i nie istnieje
ujemny cykl leżący na spacerze z $s$ do $t$,
to $\delta(s,t)$ jest kosztem pewnej
ścieżki bez powtórzeń wierzchołków.
\end{lemat}
\begin{proof}
Każdy zamknięty fragment spaceru ma wagę
nieujemną, więc możemy go usunąć bez
zwiększania kosztu. Po skończenie wielu
usunięciach otrzymamy ścieżkę prostą.
Ponieważ graf ma skończenie wiele ścieżek
prostych, najmniejszy koszt jest osiągany.
\end{proof}

\begin{definicja}[Relaksacja łuku]
\label{def:tg-weighted-relax}
Mając wartości tymczasowe $d[v]$,
dla łuku $u\to v$ wykonujemy
$d[v]\gets\min\{d[v],d[u]+w(u,v)\}$;
gdy następuje ścisła poprawa, ustawiamy
rodzica $p[v]=u$.
Początkowo $d[s]=0$, a dla $v\ne s$
przyjmujemy $d[v]=+\infty$.
\end{definicja}

\begin{lemat}[Niezmiennik relaksacji]
\label{lem:tg-weighted-walk}
Każda skończona wartość $d[v]$
otrzymana z powyższej inicjalizacji
przez relaksacje jest kosztem
rzeczywistego spaceru z $s$ do $v$.
\end{lemat}
\begin{proof}
Indukcja po liczbie relaksacji.
Wartość $d[s]=0$ odpowiada spacerowi
pustemu. Każda nowa skończona wartość
powstaje przez przedłużenie istniejącego
spaceru do $u$ o łuk $u\to v$.
\end{proof}

\section{Algorytm Dijkstry}
\label{sec:tg-weighted-dijkstra}

Zakładamy $w(a)\geq0$ dla każdego $a\in A$.
Utrzymujemy kolejkę priorytetową z kluczem $d[v]$.
Pobieramy wierzchołek o najmniejszym kluczu,
ustalamy jego odległość na stałe i relaksujemy
wszystkie wychodzące łuki. W implementacji
z powielanymi wpisami w kopcu pomijamy
wpisy nieaktualne i wierzchołki już ustalone.

\begin{twierdzenie}[Poprawność Dijkstry]
\label{thm:tg-weighted-dijkstra}
Dla skończonego grafu skierowanego
o nieujemnych wagach po ustaleniu
wierzchołka $v$ zachodzi
$d[v]=\delta(s,v)$. Po zakończeniu
wszystkie odległości są prawidłowe,
a wierzchołki nieosiągalne mają
wartość $+\infty$.
\end{twierdzenie}
\begin{proof}
Indukcja po liczbie ustalonych wierzchołków.
Początkowo $s$ ma poprawną odległość zero.
Gdyby kolejny ustalany wierzchołek $v$
miał $d[v]>\delta(s,v)$, wybierzmy
najkrótszą ścieżkę z $s$ do $v$
i jej pierwszy nieustalony wierzchołek $y$.
Jego poprzednik $x$ jest już ustalony.
Po relaksacji łuku $x\to y$ otrzymujemy
$d[y]\leq\delta(s,x)+w(x,y)=\delta(s,y)$.
Lemat \ref{lem:tg-weighted-walk} daje
równość. Nieujemne wagi dalszej części
ścieżki prowadzą do
$d[y]\leq\delta(s,v)<d[v]$,
sprzecznie z wyborem najmniejszego
klucza. Nieosiągalny wierzchołek
nie może być odkryty przez relaksację;
każdy osiągalny zostanie odkryty
po przetworzeniu poprzednika na
pewnej ścieżce z $s$.
\end{proof}

\begin{stwierdzenie}[Koszt Dijkstry]
\label{prop:tg-weighted-dijkstra-cost}
Dla list sąsiedztwa i kopca binarnego
z powielanymi wpisami otrzymujemy
$O(n+(n+m)\log(n+m+1))$ czasu
i $O(n+m)$ pamięci pomocniczej.
Jeśli kopiec obsługuje zmniejszanie klucza
z jednym wpisem na wierzchołek,
czas jest $O((n+m)\log(n+1))$,
a pamięć pomocnicza $O(n)$.
\end{stwierdzenie}
\begin{proof}
Każdy wierzchołek ustalamy co najwyżej raz,
a każdy łuk relaksujemy co najwyżej raz.
W pierwszym wariancie może powstać $O(n+m)$
wpisów kopca, a w drugim najwyżej $n$.
Złożoność operacji kopca jest
logarytmiczna względem jego rozmiaru.
\end{proof}

\begin{przyklad}[Koszt a liczba krawędzi]
\label{ex:tg-weighted-sample}
Łuki $s\to a$, $s\to b$, $a\to b$,
$a\to t$, $b\to t$ mają odpowiednio
wagi $2,5,1,7,2$.
Najtańsza droga $s\to a\to b\to t$
ma koszt $5$. Krótsza pod względem
liczby łuków droga $s\to b\to t$
ma większy koszt $7$.
\end{przyklad}

\begin{figure}[htbp]
\centering
\begin{tikzpicture}[>=Stealth,scale=.9]
\node[tgBlueVertex] (s) at (0,0) {$s$};
\node[tgVertex] (a) at (2,1.4) {$a$};
\node[tgVertex] (b) at (4,0) {$b$};
\node[tgTealVertex] (t) at (6,0) {$t$};
\draw[tgTree,->] (s)--node[above left,font=\small] {$2$} (a);
\draw[tgTree,->] (a)--node[above right,font=\small] {$1$} (b);
\draw[tgTree,->] (b)--node[above,font=\small] {$2$} (t);
\draw[tgGhost,->] (s)--node[below,font=\small] {$5$} (b);
\draw[tgGhost,->] (a) to[bend left=30] node[above,font=\small] {$7$} (t);
\end{tikzpicture}
\caption{Najtańsza droga ważona ma koszt $5$,
chociaż istnieje droga o mniejszej liczbie łuków.}
\label{fig:tg-weighted-sample}
\end{figure}

\section{Bellman--Ford i cykle o ujemnej wadze}
\label{sec:tg-weighted-bellman}

Dopuszczamy wagi ujemne. Po inicjalizacji
relaksujemy wszystkie łuki $n-1$ razy,
a następnie wykonujemy dodatkową kontrolę.

\begin{twierdzenie}[Bellman--Ford]
\label{thm:tg-weighted-bellman}
Jeżeli graf nie ma cyklu ujemnego
osiągalnego z $s$, to po $n-1$
pełnych przebiegach relaksacji
otrzymujemy $d[v]=\delta(s,v)$.
Jeśli dodatkowy przebieg pozwala
zmniejszyć oszacowanie na
osiągalnym łuku, to istnieje
ujemny cykl osiągalny z $s$,
i odwrotnie. Czas wynosi
$O(nm+n)$, pamięć pomocnicza $O(n)$.
\end{twierdzenie}
\begin{proof}
Po $k$ przebiegach dla każdej ścieżki
o co najwyżej $k$ łukach jej koszt
jest nie mniejszy od bieżącego $d$
na końcu ścieżki: dowód przez indukcję
po $k$ i relaksację ostatniego łuku.
Przy braku ujemnych cykli wystarczy
uwzględnić ścieżki proste z co
najwyżej $n-1$ łukami.
Z lematu \ref{lem:tg-weighted-walk}
wynika wtedy równość.

Jeżeli istnieje osiągalny cykl ujemny,
po $n-1$ przebiegach jego wierzchołki
mają skończone oszacowania.
Gdyby żaden łuk tego cyklu nie
umożliwiał poprawy, zsumowanie
nierówności $d[v]\leq d[u]+w(u,v)$
wzdłuż cyklu dałoby
$0\leq w(C)<0$, sprzeczność.
W przypadku braku cyklu ujemnego
wszystkie wartości są już optymalne,
więc relaksacja nie może ich poprawić.
\end{proof}

\begin{uwaga}[Odległości równe minus nieskończoności]
Jeżeli istnieje osiągalny cykl ujemny,
$-\infty$ otrzymują dokładnie
te wierzchołki, do których można
dojść z osiągalnego cyklu ujemnego.
Nie oznacza to, że każda odległość
od $s$ jest równa $-\infty$.
\end{uwaga}

\begin{przyklad}
Dla łuków $s\to a$ o wadze $3$,
$s\to b$ o wadze $7$, $a\to b$
o wadze $-5$ i $b\to t$ o wadze $2$
mamy $\delta(s,b)=-2$ i $\delta(s,t)=0$.
Dijkstra nie spełnia tutaj swoich założeń.
\end{przyklad}

\section{Floyd--Warshall: wszystkie pary}
\label{sec:tg-weighted-floyd}

Przy numeracji wierzchołków $1,\ldots,n$
ustawiamy $d^{(0)}_{ii}=0$ i
$d^{(0)}_{ij}=w(i,j)$ dla istniejących łuków,
a w pozostałych miejscach $+\infty$.
Przy łukach równoległych bierzemy minimum.
Dla $k=1,\ldots,n$ wykonujemy
\[
d^{(k)}_{ij}=
\min\{d^{(k-1)}_{ij},
d^{(k-1)}_{ik}+d^{(k-1)}_{kj}\}.
\]

\begin{twierdzenie}[Floyd--Warshall]
\label{thm:tg-weighted-floyd}
Jeżeli nie ma cykli ujemnych,
to $d^{(n)}_{ij}=\delta(i,j)$.
Graf ma cykl ujemny wtedy i tylko
wtedy, gdy wynikowa macierz ma
ujemny element przekątnej.
Koszt wynosi $O(n^3)$ czasu
i $O(n^2)$ pamięci.
\end{twierdzenie}
\begin{proof}
Indukcja po $k$: wartość $d^{(k)}_{ij}$
jest minimalnym kosztem spaceru
z $i$ do $j$ z wewnętrznymi
wierzchołkami w $\{1,\ldots,k\}$,
przy braku ujemnych cykli.
Ścieżka optymalna albo pomija $k$,
albo rozcina się w $k$ na dwa
fragmenty, co daje rekurencję.
Ujemny element przekątnej
odpowiada ujemnemu zamkniętemu
spacerowi, zatem zawiera ujemny cykl.
Z drugiej strony obecność ujemnego
cyklu daje ujemną wartość
na przekątnej po dopuszczeniu
wszystkich wierzchołków.
\end{proof}

\section{Algorytm A* i poprawna heurystyka}
\label{sec:tg-weighted-astar}

Niech $w(a)\geq0$ i niech
$t$ będzie ustalonym celem.
Heurystyka $h\colon V\to\mathbb R_{\geq0}$
z $h(t)=0$ jest \emph{dopuszczalna},
gdy $h(v)\leq\delta(v,t)$
dla wszystkich $v$ mogących
dojść do $t$. Jest \emph{spójna},
gdy $h(u)\leq w(u,v)+h(v)$
dla każdego łuku $u\to v$.
Spójność implikuje dopuszczalność.

Algorytm A* wybiera wierzchołki
z kolejki według
$f(v)=g(v)+h(v)$, gdzie
$g(v)$ jest najlepszym dotąd
kosztem od $s$ do $v$.
Przy poprawie $g(v)$
aktualizujemy też rodzica.
Dla heurystyki spójnej
możemy trwale zamknąć
wierzchołek po pierwszym
pobraniu z aktualnym kluczem.

\begin{twierdzenie}[Poprawność A*]
\label{thm:tg-weighted-astar}
Niech graf skierowany ma
nieujemne wagi, a $h$
będzie nieujemna i spójna,
przy czym $h(t)=0$.
Algorytm A* bez ponownego
otwierania zamkniętych
wierzchołków zwraca
$\delta(s,t)$ albo
stwierdza brak drogi.
\end{twierdzenie}
\begin{proof}
Wprowadźmy zredukowane
wagi $w_h(u,v)=w(u,v)+h(v)-h(u)$,
które są nieujemne ze spójności.
Wzdłuż dowolnego spaceru
$P$ od $s$ do $v$ zachodzi
$w_h(P)=w(P)+h(v)-h(s)$.
Zatem najkrótsze drogi nie
zmieniają się po redukcji.
Klucz Dijkstry dla wag
$w_h$ jest równy
$g(v)+h(v)-h(s)$,
czyli kluczowi A*
pomniejszonemu o stałą $h(s)$.
Algorytmy wybierają więc
wierzchołki w tej samej
kolejności, a poprawność
wynika z twierdzenia
\ref{thm:tg-weighted-dijkstra}.
\end{proof}

\begin{uwaga}[Ponowne otwieranie]
Jeżeli heurystyka jest
dopuszczalna, ale niespójna,
poprawa kosztu $g(v)$
może nadejść po zamknięciu $v$.
Dla zachowania dokładności
trzeba wtedy ponownie
otwierać wierzchołki.
Nie wolno automatycznie
stosować wariantu
bez ponownych otwarć.
Dla $h\equiv0$ A*
redukuje się do Dijkstry.
\end{uwaga}

\begin{przyklad}[Heurystyka euklidesowa]
Jeżeli wierzchołki mają
pozycje $x_v\in\mathbb R^2$
i $w(u,v)\geq\|x_u-x_v\|_2$,
to $h(v)=\|x_v-x_t\|_2$
jest spójna na mocy
nierówności trójkąta.
Przy koszcie czasu
jazdy, a nie długości,
potrzebne jest dodatkowe
przeskalowanie uwzględniające
górne ograniczenie prędkości.
\end{przyklad}

\section{Dwukierunkowy algorytm Dijkstry}
\label{sec:tg-weighted-bidir}

Dla zapytania $s\to t$
prowadzimy Dijkstrę
od $s$ po zwykłych
łukach oraz od $t$
po łukach odwróconych.
Niech $d_F(v)$ i $d_B(v)$
oznaczają znane koszty
dojścia do $v$ odpowiednio
z $s$ i wstecz z $t$.
Gdy obie wartości są
skończone, koszt
$d_F(v)+d_B(v)$
jest górnym ograniczeniem
odległości od $s$ do $t$.
Przechowujemy minimum
tych kandydatów i
kandydatów przez łączące
fronty łuki, oznaczone $\mu$.

\begin{stwierdzenie}[Kryterium zakończenia]
\label{prop:tg-weighted-bidir}
Dla grafu o nieujemnych
wagach, gdy minima aktualnych
kluczy w dwóch kolejkach
są równe $k_F,k_B$,
warunek $k_F+k_B\geq\mu$
pozwala bezpiecznie
zakończyć wyszukiwanie
i zwrócić $\mu$.
Samo pierwsze spotkanie
frontów nie jest wystarczające.
\end{stwierdzenie}
\begin{proof}
Każda skończona wartość
$\mu$ pochodzi z
rzeczywistego spaceru,
więc nie jest mniejsza
od $\delta(s,t)$.
Jeśli istniałaby tańsza
ścieżka $P$, na jej
granicy z obszarami
ustalonymi z obu stron
znaleźlibyśmy albo
wierzchołek, albo
łuk łączący fronty
i dający kandydata
o koszcie co najwyżej
$w(P)$. Dopóki taki
łącznik nie jest dostępny,
nieustalone fragmenty
tej ścieżki mają
łączny koszt co najmniej
$k_F+k_B$.
Przy $k_F+k_B\geq\mu$
prowadzi to do sprzeczności.
\end{proof}

\section{Hierarchie kontrakcji: szybkie zapytania po przygotowaniu grafu}
\label{sec:tg-weighted-ch}

Opisujemy dokładną wersję
dla skończonego grafu
\emph{nieskierowanego}
o \emph{dodatnich} wagach.
Ustalmy bijekcję rang
$\rho\colon V\to\{1,\ldots,n\}$.
Kontraktujemy kolejne
wierzchołki od
najniższej rangi.
Przy usuwaniu $v$,
dla każdej pary jego
aktualnych sąsiadów
$u,z$ o wyższej
randze dodajemy
\emph{skrót} $uz$
o wadze $w(u,v)+w(v,z)$,
jeśli nie ma jeszcze
krawędzi o wadze
nie większej.
Każdy skrót pamięta,
jaką parę krawędzi
zastąpił, co pozwala
odtworzyć trasę.

\begin{lemat}[Zachowanie odległości]
\label{lem:tg-weighted-ch}
Po kontrakcji wierzchołków
o rangach $1,\ldots,k$
odległości między
pozostałymi wierzchołkami
w grafie aktywnym są
równe odległościom
w grafie oryginalnym.
\end{lemat}
\begin{proof}
Indukcja po $k$.
Po usunięciu $v$
każdy fragment
$u-v-z$ można zastąpić
skrótem o nie większej
wadze, więc odległości
nie wzrastają.
Każdy nowy skrót
odpowiada zarazem
spacerowi w poprzednim
grafie tej samej wagi,
więc odległości
nie maleją.
\end{proof}

W grafie rozszerzonym
zachowujemy wszystkie
oryginalne krawędzie
i dodane skróty,
ale wyszukiwanie
prowadzimy tylko
po łukach o
\emph{rosnącej} randze,
czyli w grafie
$D_\uparrow$.
Niech $d_\uparrow(u,v)$
oznacza odległość
w takim grafie
skierowanym.

\begin{twierdzenie}[Dokładne zapytanie CH]
\label{thm:tg-weighted-ch-query}
Dla dowolnych wierzchołków
$s,t$ oryginalnego
grafu zachodzi
\[
\delta(s,t)=
\min_{x\in V}
\{d_\uparrow(s,x)+d_\uparrow(t,x)\}.
\]
Dwie Dijkstry po
łukach rosnącej rangi,
z $s$ oraz z $t$,
pozwalają obliczyć
prawą stronę.
\end{twierdzenie}
\begin{proof}
Połączenie obu ścieżek
w górę w punkcie $x$
daje spacer od $s$ do
$t$ w grafie rozszerzonym.
Rozwinięcie skrótów
nie zmienia kosztu,
więc prawa strona
nie jest mniejsza
od $\delta(s,t)$.

Weźmy ścieżkę najkrótszą.
Jeśli jej ciąg rang
ma lokalne minimum,
to podczas kontrakcji
odpowiedniego wierzchołka
utworzono skrót między
jego sąsiadami albo
pozostawiono krawędź
nie droższą.
Zastąpmy ten fragment.
Usuwając kolejno minima
o najmniejszej randze,
po skończonej liczbie
kroków otrzymamy
spacer, którego rangi
najpierw rosną, a
potem maleją.
Obie połowy od końców
do wierzchołka
o największej randze
są ścieżkami w górę.
Prawa strona nie
przekracza więc
$\delta(s,t)$.
\end{proof}

\begin{uwaga}[Koszt przygotowania]
Wariant dodający wszystkie
skróty jest dokładny, ale
w najgorszym przypadku
może stworzyć graf gęsty.
Praktyczne CH ograniczają
liczbę skrótów za pomocą
poszukiwania ścieżek
omijających kontraktowany
wierzchołek i dobrze
dobranej kolejności rang.
Zmiana wag może wymagać
przebudowy struktury.
Nie jest to zatem
uniwersalnie najszybsza
metoda dla pojedynczego
zapytania; jej siła
tkwi często w wielu
zapytaniach na
tej samej sieci drogowej.
\end{uwaga}

\section{Porównanie i wybór algorytmu}
\label{sec:tg-weighted-summary}

\begin{center}
\small
\begin{tabular}{p{.23\textwidth}p{.29\textwidth}p{.33\textwidth}}
\hline
Algorytm & Założenia & Zastosowanie\\
\hline
BFS & koszty jednostkowe & graf nieważony\\
Dijkstra & wagi nieujemne & jedno źródło lub cel\\
Bellman--Ford & dopuszczalne wagi ujemne & jedno źródło, wykrywanie cykli\\
Floyd--Warshall & graf skończony & wszystkie pary\\
A* & nieujemne wagi, dobra heurystyka & pojedyncza trasa\\
Dwukierunkowy Dijkstra & wagi nieujemne & pojedyncza para\\
CH & dodatnie wagi, przygotowanie grafu & wiele tras w stabilnej sieci\\
\hline
\end{tabular}
\end{center}

A* zwykle przetwarza mniej
wierzchołków, gdy potrafimy
podać informacyjną,
dopuszczalną heurystykę.
CH może oferować
znacznie krótsze czasy
zapytań na sieciach
drogowych po kosztownym
przygotowaniu indeksu.
Każde porównanie
wymaga rzeczywistych
danych i benchmarku.

\begin{uwaga}[Literatura]
R. Geisberger, P. Sanders,
D. Schultes, D. Delling,
\emph{Contraction Hierarchies:
Faster and Simpler Hierarchical
Routing in Road Networks},
WEA 2008.
\end{uwaga}
